% Figure: photoelectric stopping voltage versus frequency. The maximum
% kinetic energy of emitted electrons rises linearly with photon frequency
% above the threshold; the slope is Planck's constant and the intercept is
% the work function. Two metals shown to make the parallel-line / shifted
% threshold point.
\begin{figure}[htbp]
\centering
\begin{tikzpicture}
\begin{axis}[
  width=12cm, height=7.5cm,
  xlabel={photon frequency $\nu$ (PHz)},
  ylabel={stopping voltage $V_0$ (V)},
  xmin=0, xmax=1.6, ymin=-0.5, ymax=5.2,
  axis lines=left,
  legend pos=north west,
  legend cell align=left,
  tick label style={font=\scriptsize},
  label style={font=\small},
  legend style={font=\scriptsize},
  clip=false,
]
% cesium: work function ~2.1 eV, threshold ~0.51 PHz. slope h/e ~4.136 eV/PHz.
\addplot[engblue, very thick, domain=0.508:1.55, samples=2] {4.136*x - 2.1};
\addlegendentry{cesium, $\phi \approx 2.1$ eV}
% silver: work function ~4.7 eV, threshold ~1.14 PHz.
\addplot[engred, very thick, domain=1.136:1.55, samples=2] {4.136*x - 4.7};
\addlegendentry{silver, $\phi \approx 4.7$ eV}
% threshold markers
\node[engblue, font=\scriptsize, anchor=north] at (axis cs:0.508,-0.15) {$\nu_{0,\text{Cs}}$};
\node[engred, font=\scriptsize, anchor=north] at (axis cs:1.136,-0.15) {$\nu_{0,\text{Ag}}$};
\node[font=\scriptsize, anchor=west] at (axis cs:0.9,3.6) {slope $= h/e$};
\end{axis}
\end{tikzpicture}
\caption[Photoelectric stopping voltage versus frequency]{Stopping voltage
$V_0$ as a function of incident photon frequency, for two photocathode
metals. Below a threshold frequency $\nu_0$ no current flows at any
intensity. Above threshold the maximum electron kinetic energy is
$eV_0 = h\nu - \phi$, so the data fall on a straight line of slope $h/e$
that is the same for every material; only the intercept, set by the work
function $\phi$, shifts. The common slope is the measurement that fixed
Planck's constant from the photoelectric effect and confirmed that light
delivers energy in quanta of size $h\nu$.}
\label{fig:vol04ch12:photoelectric-stopping}
\end{figure}
